% Auto-generated by src/analysis/latex_tables.py — do not edit by hand.
\begin{table}[htbp]
\centering
\small
\caption{Princeton fertility decomposition (Coale and Watkins 1986): marital fertility $I_g$ vs.\ nuptiality (marriage rate)}
\label{tab:coale_decomposition}
\begin{tabular}{lcccc}
\toprule
 & \multicolumn{2}{c}{Low Catholic ($\le$50\%)} & \multicolumn{2}{c}{High Catholic ($>$50\%)} \\
\cmidrule(lr){2-3}\cmidrule(lr){4-5}
 & Pre & Post & Pre & Post \\
\midrule
\multicolumn{5}{l}{\textit{Panel A: Group means}} \\
$I_f$ (overall fertility) & 0.408 & 0.420 & 0.425 & 0.437 \\
$I_g$ (marital fertility) & 0.706 & 0.712 & 0.795 & 0.827 \\
$I_h$ (illegitimate fertility) & 0.079 & 0.080 & 0.048 & 0.045 \\
Marriage rate (per 1{,}000 pop) & 8.405 & 8.034 & 8.101 & 7.484 \\
Gen.\ marriage rate (per 1{,}000 pop 15+) & 13.164 & 12.561 & 13.013 & 11.946 \\
\midrule
\multicolumn{5}{l}{\textit{Panel B: DiD coefficient on $\mathrm{CathShare} \times \mathrm{Post}$ (coefficient, SE, observations)}} \\
Outcome & Coefficient & SE & $N$ & \\
\cmidrule(lr){1-4}
$I_f$ (overall fertility) & -0.00007 & (0.00006) & 10,783 \\
$I_g$ (marital fertility) & +0.00027$^{**}$ & (0.00011) & 10,783 \\
$I_h$ (illegitimate fertility) & -0.00010$^{***}$ & (0.00002) & 10,783 \\
Marriage rate (per 1{,}000 pop) & -0.004$^{***}$ & (0.001) & 10,777 \\
Gen.\ marriage rate (per 1{,}000 pop 15+) & -0.008$^{***}$ & (0.002) & 10,783 \\
\bottomrule
\end{tabular}
\begin{minipage}{\linewidth}
\vspace{0.5em}
\footnotesize \textit{Notes:} $I_f$, $I_g$, and $I_h$ are Hutterite-benchmarked Princeton EFP indices for overall, marital, and non-marital fertility. They are computed using the Coale-Demeny ``West'' female age distribution and an assumed nuptiality schedule (calibrated to a weighted-mean marriage prevalence of 64\% among women 15--49, consistent with the eastern side of the Hajnal line). Because Galloway lacks age structure of women and married women, the absolute index levels depend on the calibration and should be read alongside the cross-county and pre/post differences. The Princeton I_m index is omitted because it is constant under the calibration; the observed marriage rate (Martot / Poptot per 1{,}000) serves as the nuptiality marker. Panel~A reports group means; Panel~B reports the coefficient on $\mathrm{CathShare} \times \mathrm{Post}$ from a county- and year-fixed-effects regression with $\ln(\mathrm{Pop})$ as a control. The demographic interpretation: $I_g$ shows no DiD effect, but the marriage rate shows a strong negative effect, implying that the Kulturkampf operated through marriage formation (nuptiality), not within-marriage childbearing. Standard errors clustered at the county level. $^{*}\,p<0.10$, $^{**}\,p<0.05$, $^{***}\,p<0.01$.
\end{minipage}
\end{table}
